\chapter{Antecedentes}
\label{cap:antecedentes}

\section{Evolución de la gestión de reservas y paquetes turísticos}

Antes de la expansión de Internet, muchas agencias de turismo atendían las
consultas directamente en sus oficinas o por teléfono. Los catálogos
impresos y los contactos con empresas de transporte y hospedaje servían
para presentar las ofertas. Con estos medios era posible trabajar, siempre
que se mantuviera un control adecuado de los cambios y las confirmaciones.

Posteriormente se extendió el uso de sistemas computarizados de reservas.
Los sistemas globales de distribución, conocidos como GDS, facilitaron la
consulta de horarios, tarifas y disponibilidad de diferentes proveedores
turísticos desde un punto de acceso común \parencite{amadeusGDS}. Esta
evolución permitió organizar mejor la información de servicios que antes se
consultaban por separado.

Con el crecimiento de los servicios en línea, los clientes comenzaron a
buscar y reservar opciones de viaje desde sus computadoras y celulares.
Booking.com, cuya actividad comenzó en 1996, es un ejemplo de la expansión
de las plataformas digitales en el sector turístico \parencite{booking}.
Sin embargo, una agencia de Potosí que ofrece paquetes propios no necesita
copiar todas las funciones de una plataforma internacional. Lo principal es
que pueda administrar correctamente sus servicios y cupos.

En los últimos años se han desarrollado herramientas para elaborar
cotizaciones, organizar itinerarios, guardar información de proveedores y
controlar pagos. Por ejemplo, Ezus presenta una solución que reúne
cotizaciones, itinerarios, tarifas y documentación \parencite{ezus}. Estas
funciones muestran cómo distintos procesos se pueden relacionar dentro de un
mismo sistema, aunque cada agencia debe seleccionar lo que realmente
necesita.

Para este proyecto, la variable que se busca mejorar es la consistencia del
proceso de reservas y del control de disponibilidad. Esta puede observarse
en la exactitud de los cupos, el tiempo que demora una confirmación y la
cantidad de errores detectados cuando se registra o modifica una reserva.

\section{Relación con sistemas similares}

Como referencia internacional se puede mencionar a Amadeus, que facilita el
acceso a información de servicios turísticos de distintos proveedores, y a
Booking.com, que permite consultar y reservar opciones de viaje en línea.
Por otra parte, programas de gestión como Ezus incluyen herramientas
administrativas para preparar propuestas, organizar itinerarios y manejar la
información comercial de una agencia \parencite{amadeusGDS,booking,ezus}.

En la región existen trabajos más cercanos a este proyecto.
\textcite{huaman2017} desarrolló un sistema de información para la gestión
de reservas de paquetes turísticos en la agencia Rap Travel, tras
identificar que el proceso se realizaba manualmente y sin un control
adecuado del seguimiento de los pagos. Por su parte,
\textcite{miranda2022} elaboraron un sistema informático con aplicación
móvil para gestionar reservas de paquetes turísticos y asesoría migratoria
en una agencia de viajes del municipio de San Vicente. A nivel nacional, en
la búsqueda realizada para este trabajo no se encontró documentación
pública sobre sistemas de gestión de reservas para agencias de turismo de
Potosí, por lo que el diagnóstico en la agencia seleccionada permitirá
documentar su situación.

Nuestro proyecto comparte con estos sistemas la idea de mantener la
información organizada y actualizada, pero tendrá un alcance menor. No se
pretende crear una plataforma mundial ni conectarse directamente a
aerolíneas: el sistema se enfocará en los paquetes que ofrece una agencia de
turismo de Potosí y en la forma de consultar, reservar y controlar los
cupos de cada salida programada. También se tomarán en cuenta criterios
básicos de usabilidad, pues la norma ISO 9241-210:2019 destaca la
importancia de diseñar sistemas interactivos considerando las necesidades
de las personas usuarias \parencite{iso9241}.

\section{Flujo de información de la agencia}

El proceso comienza cuando el cliente solicita información sobre un destino
o paquete turístico. El encargado revisa el itinerario, el precio, los
servicios incluidos y la fecha de salida. Después verifica los cupos
disponibles y prepara una cotización o una respuesta con las condiciones
del viaje. Si el cliente acepta la oferta, se registran sus datos y la
cantidad de pasajeros.

En la siguiente etapa, el personal comprueba que la reserva cumpla las
condiciones de la agencia y cambia su estado a pendiente o confirmada. Si
corresponde realizar un anticipo, se registra el pago o el comprobante
recibido. Cuando se modifica la fecha, el número de pasajeros o se cancela
una reserva, el sistema debe volver a calcular los cupos disponibles para
evitar diferencias en la información.

De manera complementaria, los datos de las reservas confirmadas pueden
utilizarse para coordinar el transporte, los guías y otros servicios
necesarios. Esta descripción representa el flujo que se propone analizar.
Su funcionamiento real se verificará en la agencia seleccionada, porque cada
empresa puede tener procedimientos distintos.
